\documentclass[14pt,a4paper,oneside]{extbook}
\usepackage[utf8]{vietnam}
\usepackage{amsmath,amssymb,mathrsfs,fancyhdr,enumerate,multirow,makecell}
\usepackage{tkz-euclide,tikz-3dplot,tikz,tkz-tab}
\usetikzlibrary{arrows.meta,calc,intersections,angles,quotes,backgrounds,shapes.geometric,patterns,through,shadows,positioning}
\usepackage{yhmath}
\usepackage[top=0.7cm, bottom=0.7cm, left=1.0cm, right=1.0cm, headheight=16pt, headsep=3mm, footskip=5mm]{geometry}
\usepackage[hidelinks,unicode]{hyperref}
\usepackage{qrcode}
\usepackage[many]{tcolorbox}
\usepackage{fontawesome}

\pagestyle{fancy}
\fancyhf{}
\fancyhead[L]{\footnotesize\bfseries\color{teal!85!black}\faBook~SỔ TAY KIẾN THỨC TOÁN 4}
\fancyhead[R]{\footnotesize\bfseries\color{blue!85!black}Lớp Toán Thầy Lĩnh -- loptoanthaylinh.vn}
\fancyfoot[C]{\footnotesize\bfseries\color{gray!80!black}-- Sổ tay kiến thức: Bài 66. Luyện tập chung nhân chia phân số --}
\renewcommand{\headrulewidth}{0.4pt}

\begin{document}

% --- BANNER TIÊU ĐỀ BÀI HỌC ---
\noindent
\begin{tcolorbox}[
  enhanced,
  colback=teal!8!white,
  colframe=teal!80!black,
  arc=3mm, boxrule=1pt,
  left=3mm, right=3mm, top=2mm, bottom=2mm,
  halign=center
]
  {\Large\bfseries\color{teal!90!black}\faBookmark\ BÀI 66: LUYỆN TẬP CHUNG NHÂN CHIA PHÂN SỐ}\\[1.5pt]
  {\footnotesize\itshape\color{black!75} Thứ tự thực hiện phép tính, tính chất phân phối và tìm thành phần chưa biết trên 1 trang A4}
\end{tcolorbox}

\vspace{-2.5mm}

\noindent
\begin{minipage}[t]{0.49\textwidth}\vspace{0pt}
\begin{tcolorbox}[
  enhanced,
  colback=blue!4!white,
  colframe=blue!75!black,
  arc=2mm, boxrule=0.8pt,
  left=2mm, right=2mm, top=1.8mm, bottom=1.8mm,
  title={\footnotesize\bfseries\color{white}\faListOl\ 1. Thứ tự thực hiện phép tính}
]
{\footnotesize
Trong biểu thức có phân số:
\begin{itemize}\setlength{\itemsep}{1pt}\setlength{\leftskip}{-2.5mm}
  \item \textbf{Có ngoặc:}\\
  Thực hiện trong ngoặc $(~)$ trước, ngoài ngoặc sau.
  \item \textbf{Không có ngoặc:}\\
  Thực hiện \textbf{Nhân, Chia trước}; \textbf{Cộng, Trừ sau}.
  \item \textbf{Chỉ có Nhân và Chia (hoặc chỉ Cộng và Trừ):}\\
  Thực hiện lần lượt từ \textbf{trái sang phải}.
\end{itemize}
}
\end{tcolorbox}
\end{minipage}
\hfill
\begin{minipage}[t]{0.49\textwidth}\vspace{0pt}
\begin{tcolorbox}[
  enhanced,
  colback=teal!4!white,
  colframe=teal!75!black,
  arc=2mm, boxrule=0.8pt,
  left=2mm, right=2mm, top=1.8mm, bottom=1.8mm,
  title={\footnotesize\bfseries\color{white}\faMagic\ 2. Tính chất phân phối nhân với cộng}
]
{\footnotesize
$$\dfrac{a}{b} \times \dfrac{c}{d} + \dfrac{a}{b} \times \dfrac{p}{q} = \dfrac{a}{b} \times \left(\dfrac{c}{d} + \dfrac{p}{q}\right)$$
$$\dfrac{a}{b} \times \dfrac{c}{d} - \dfrac{a}{b} \times \dfrac{p}{q} = \dfrac{a}{b} \times \left(\dfrac{c}{d} - \dfrac{p}{q}\right)$$
\textbf{Bí quyết:} Đặt thừa số chung $\dfrac{a}{b}$ ra ngoài để tính trong ngoặc ra số nguyên đẹp!
}
\end{tcolorbox}
\end{minipage}

\vspace{-2.5mm}
% --- KHỐI 3: SƠ ĐỒ TƯ DUY TIKZ ---
\noindent
\begin{tcolorbox}[
  enhanced,
  colback=orange!4!white,
  colframe=orange!80!black,
  arc=2mm, boxrule=0.8pt,
  left=2mm, right=2mm, top=1.5mm, bottom=1.5mm,
  title={\footnotesize\bfseries\color{white}\faSitemap\ SƠ ĐỒ TƯ DUY: HỆ THỐNG 4 KĨ NĂNG NHÂN CHIA PHÂN SỐ}
]
\centering
\begin{tikzpicture}[
  box/.style={draw, rounded corners=2pt, align=center, font=\scriptsize\bfseries, inner sep=3pt, minimum width=3.4cm, minimum height=0.9cm},
  arr/.style={-{Stealth[length=2mm]}, thick}
]
  \node[draw=purple!80!black, fill=purple!15, rounded corners=3pt, font=\small\bfseries\color{purple!90!black}, inner sep=4pt, align=center] (core) at (5, 0.7) {NHÂN CHIA\\PHÂN SỐ};

  \node[box, fill=blue!10, draw=blue!80!black] (k1) at (0.8, 1.4) {1. PHÉP NHÂN\\Tử $\times$ Tử, Mẫu $\times$ Mẫu};
  \node[box, fill=teal!10, draw=teal!80!black] (k2) at (9.2, 1.4) {2. PHÉP CHIA\\Nhân đảo ngược số chia};
  \node[box, fill=orange!12, draw=orange!80!black] (k3) at (0.8, 0.0) {3. TÌM PHÂN SỐ\\Lấy số nhân phân số};
  \node[box, fill=green!15, draw=green!70!black] (k4) at (9.2, 0.0) {4. PHÂN PHỐI\\Đặt thừa số chung};

  \draw[arr, blue!80!black] (core) -- (k1);
  \draw[arr, teal!80!black] (core) -- (k2);
  \draw[arr, orange!80!black] (core) -- (k3);
  \draw[arr, green!70!black] (core) -- (k4);
\end{tikzpicture}
\end{tcolorbox}

\vspace{-2.5mm}
% --- KHỐI 4: VÍ DỤ & MẸO TRÁNH BẪY ---
\noindent
\begin{minipage}[t]{0.49\textwidth}\vspace{0pt}
\begin{tcolorbox}[
  enhanced,
  colback=green!4!white,
  colframe=green!60!black,
  arc=2mm, boxrule=0.8pt,
  left=2mm, right=2mm, top=1.8mm, bottom=1.8mm,
  title={\footnotesize\bfseries\color{white}\faCheckCircle\ Ví dụ thực hành chuẩn mực}
]
{\footnotesize
\textbf{1. Tính thuận tiện (Đặt thừa số chung):}\\
$B = \dfrac{3}{5} \times \dfrac{4}{7} + \dfrac{3}{5} \times \dfrac{3}{7} = \dfrac{3}{5} \times \left(\dfrac{4}{7} + \dfrac{3}{7}\right)$\\
$B = \dfrac{3}{5} \times \dfrac{7}{7} = \dfrac{3}{5} \times 1 = \mathbf{\dfrac{3}{5}}$.\\
\textbf{2. Tìm $x$:}\\
$x \times \dfrac{2}{3} = \dfrac{4}{9} \implies x = \dfrac{4}{9} : \dfrac{2}{3} = \dfrac{4}{9} \times \dfrac{3}{2} = \mathbf{\dfrac{2}{3}}$.
}
\end{tcolorbox}
\end{minipage}
\hfill
\begin{minipage}[t]{0.49\textwidth}\vspace{0pt}
\begin{tcolorbox}[
  enhanced,
  colback=red!4!white,
  colframe=red!70!black,
  arc=2mm, boxrule=0.8pt,
  left=2mm, right=2mm, top=1.8mm, bottom=1.8mm,
  title={\footnotesize\bfseries\color{white}\faExclamationTriangle\ Mẹo tránh bẫy trắc nghiệm}
]
{\footnotesize
\begin{itemize}\setlength{\itemsep}{1.5pt}\setlength{\leftskip}{-2.5mm}
  \item \textbf{Bẫy tìm số chia $x$:}\\
  $\frac{2}{3} : x = \frac{4}{5} \implies x = \frac{2}{3} : \frac{4}{5}$ (Số bị chia chia cho Thương). Học sinh hay nhầm thành phép nhân!
  \item \textbf{Bẫy thứ tự trái sang phải:}\\
  Biểu thức $\frac{3}{4} : \frac{1}{2} \times \frac{2}{3}$ phải thực hiện chia trước rồi mới nhân sau, không được làm phép nhân trước!
\end{itemize}
}
\end{tcolorbox}
\end{minipage}

\end{document}
